% Experimental setup. Sources: EVIDENCE.md A-E, run_*_flyoff.sh, video_streams.py
\section{Experimental Setup}
\label{sec:setup}

\textit{Traffic and metrics.} Each robot sends constant-rate UDP (User Datagram Protocol) video at
1.5\,Mb/s to the core and is probed for C2 by ping; a robot's C2 is up when
fewer than 50\,\% of its probes are lost. These definitions differ from those
of~\cite{makhubela2026satnac} (two robots, 8\,Mb/s video bursts, C2 up when
any of five echoes returns), so numbers are not comparable across the two
papers. Placement runs are reported as the time-weighted share of a mission
with mean C2 up, with all four robots up, and with no robot up, and the
90th-percentile video latency. Forecasts are scored by the area under the
precision--recall curve (PR-AUC) relative to the no-skill rate (lift) and by ROC-AUC, pooled over held-out terrains and on the
worst held-out terrain.

\textit{Data.} The forecasting and altitude-effect results use 54,572
robot-cycles from 20 missions of 90--200\,min over four terrains (open field,
urban, forest, hills), all flown before the cap existed. The cap was tested in
a static load test in forest and open field (6\,min of video between 2-min
no-load phases), in a 30-min forest patrol, and in the placement fly-offs.

\textit{Placement arms.} Every arm flies the same pilot, seed and 90-min
mission and differs only in how x--y and altitude are chosen
(Table~\ref{tab:arms}): the pilot decides every 30\,s, changes altitude by at
most 10\,m per decision, and stays within a 20--60\,m geofence. The 40-m hover altitude and the
20--60\,m band are those of~\cite{makhubela2026satnac}: 40\,m is the middle of
the band and clears every obstacle in the outdoor worlds. Arm B, the
robot centroid at fixed altitude, is the baseline of our research proposal; arm
A, fixed hover, is the simplest policy; the repositioning policy
of~\cite{makhubela2026satnac} is a further baseline for the risk planner.

\begin{table}[t]
\caption{Placement Arms}
\label{tab:arms}
\centering
\small
% Table 1 (hand-written: definitions, no numbers). Source: run_xyz_flyoff.sh,
% scripts/ml/online_planner.py (SAR_PLANNER_XY / _ALT / _ALT_LEARNER).
\begin{tabular}{@{}llll@{}}
\toprule
Arm & x--y & Altitude & Role \\
\midrule
A & held & fixed 40\,m & baseline (hover) \\
B & robot centroid & fixed 40\,m & baseline (proposal) \\
C & robot centroid & learned on the spot & learned altitude \\
D & risk model & learned on the spot & learned x--y--z \\
\bottomrule
\end{tabular}

\end{table}

\textit{Validation protocol.} A result is accepted only after it survives the
control that matches its threat: leave-one-terrain-out evaluation with
worst-fold reporting for forecasting; censoring of dead-process cycles for every
label; a same-policy null control (the same comparison applied to two nights of
one policy) for closed-loop comparisons made across nights, where a
comparison counts only if the null control's interval rules out a swing as
large as the effect; and the uplink cap
for any forest outcome. Confidence intervals are 95\,\% block bootstraps over
contiguous stretches (10-min blocks for missions), because consecutive cycles
are strongly correlated; an effect whose interval includes zero is reported as
not established. Per-arm placement comparisons use one 90-min run per arm and
within-run intervals, which do not capture run-to-run variation.
\todo{update: three seeds per arm in forest and open field from the repeat campaign}
